\DocumentMetadata{
  pdfversion=2.0,
  pdfstandard=ua-2,
  lang=en-US,
  tagging=on,
  tagging-setup={math/setup=mathml-SE,tabsorder=structure}
}
\documentclass[11pt,letterpaper]{article}
\usepackage[margin=0.68in,includefoot]{geometry}
\usepackage{newcomputermodern}
\usepackage{amsmath,amscd,mathtools}
\providecommand{\square}{\mdlgwhtsquare}
\usepackage[hidelinks]{hyperref}
\hypersetup{
  pdftitle={Analysis First-Year Exam, May 2021, Part 2},
  pdfauthor={University of Florida Department of Mathematics},
  pdfsubject={Graduate mathematics examination},
  pdfdisplaydoctitle=true,
  bookmarksopen=true
}
\setcounter{secnumdepth}{0}
\setlength{\parindent}{0pt}
\setlength{\parskip}{0.28em}
\setlength{\emergencystretch}{3em}
\setlength{\leftmargini}{2.15em}
\setlength{\leftmarginii}{2.65em}
\setlength{\leftmarginiii}{3em}
\renewcommand{\labelenumi}{\textbf{\theenumi.}}
\pagestyle{empty}
\raggedbottom
\allowdisplaybreaks
\newenvironment{examproblems}[1]{%
  \begin{enumerate}%
  \setcounter{enumi}{#1}%
  \setlength{\topsep}{0.35em}%
  \setlength{\partopsep}{0pt}%
  \setlength{\itemsep}{0.48em}%
  \setlength{\parsep}{0.18em}%
}{\end{enumerate}}
\newenvironment{examsubparts}{%
  \begin{enumerate}%
  \setlength{\topsep}{0.18em}%
  \setlength{\partopsep}{0pt}%
  \setlength{\itemsep}{0.16em}%
  \setlength{\parsep}{0.1em}%
}{\end{enumerate}}
\newenvironment{examinstructions}{%
  \par\begingroup\itshape
}{%
  \par\endgroup\vspace{0.18em}
}
\makeatletter
\renewcommand\section{\@startsection{section}{1}{0pt}%
  {-1.25ex plus -.25ex minus -.1ex}%
  {0.55ex plus .1ex}%
  {\normalfont\large\bfseries}}
\renewcommand{\@maketitle}{%
  \newpage\null\vskip -1.2em%
  {\footnotesize\bfseries
    \href{https://www.ufl.edu/}{University of Florida}, \href{https://math.ufl.edu/}{Department of Mathematics}\par}%
  \vskip 0.18em%
  \tagstructbegin{tag=Title}%
    {\LARGE\bfseries\href{https://gma.math.ufl.edu/exams/analysis-first-year/}{Analysis First-Year Exam}, May 2021, Part 2\par}%
  \tagstructend%
  \vskip 0.35em%
  \tagmcbegin{artifact}%
    \hrule height 1pt%
  \tagmcend%
  \vskip 0.55em%
}
\makeatother
\title{Analysis First-Year Exam, May 2021, Part 2}
\author{}
\date{}
\begin{document}
\maketitle
\thispagestyle{empty}
\begin{examinstructions}
Answer FOUR questions in detail. State carefully any results used without proof.
\end{examinstructions}
\begin{examproblems}{0}
\item
For each \(n\in\mathbb{N}\) let \(f_n:[0,1]\to\mathbb{R}\) be continuous. Assume that for each \(t\in[0,1]\) the sequence \((f_n(t):n\in\mathbb{N})\) converges to~\(0\) monotonically. Does it follow that \(\int_0^1 f_n(t)\,dt\to 0\)? Does the answer change if the word ‘monotonically’ is removed?
\item
Let \(f:[-1,1]\to\mathbb{R}\) be continuous and assume that \(\int_{-1}^1 f(t)t^n\,dt=0\) for each even integer \(n\geq 0\). Deduce as much as is possible about the nature of the function~\(f\).
\item
Let \((f_n:n\in\mathbb{N})\) be a sequence of measurable real-valued functions on the same measurable space. In each of the following cases, show that the set of points~\(\omega\) satisfying the stated condition is a measurable set:
\begin{examsubparts}
\item[\textnormal{(i)}]
the sequence \((f_n(\omega):n\in\mathbb{N})\) is unbounded;
\item[\textnormal{(ii)}]
\((f_n(\omega):n\in\mathbb{N})\) is not strictly monotonic;
\item[\textnormal{(iii)}]
\((f_n(\omega):n\in\mathbb{N})\) revisits its initial value \(f_0(\omega)\) infinitely often.
\end{examsubparts}
\item
Let \((\Omega,\mathcal{A},\mu)\) be a measure space on which \(f\geq 0\) is an integrable function. Prove that for each \(\varepsilon>0\) there exists \(\delta>0\) such that each \(A\in\mathcal{A}\) with \(\mu(A)<\delta\) satisfies \(\int_A f\,d\mu<\varepsilon\).
\item
Consider these two statements about the function \(f:\mathbb{R}\to\mathbb{R}\). Prove that (a)~\(\Rightarrow\) (b) and show by example that (b)~\(\Rightarrow\) (a) can fail.
\begin{examsubparts}
\item[\textnormal{(a)}]
\(f\) is integrable with respect to Lebesgue measure~\(\lambda\) on~\(\mathbb{R}\);
\item[\textnormal{(b)}]
\(f\) is Lebesgue integrable on \([-n,n]\) for each integer \(n>0\) and the sequence of integrals \(\int_{-n}^n f\,d\lambda\) converges.
\end{examsubparts}
\end{examproblems}
\end{document}
